\section{Non-perturbative results}
\label{sec:nonper}
\begin{figure}[htbp]
    \centering
    \includegraphics[width=0.6\textwidth]{images/zvl.pdf}
    \caption{\itshape $Z_{V^L}$ as a function of $\mu^2 a^2$.
Since the lattice is not symmetric in space and time, the results
obtained by using the time-component of the local current ($V_0$)
and the space components, averaged in the three directions, ($V_k$)
are shown separately. The dashed line is $Z^{{\rm BSPT}}_{V^L}$ from BSPT,
 and the curve is from BDPT,   on a
volume of size $16^3 \times 32$, see sec.\ \ref{sec:per},
see sec.\ \ref{sec:per}.
 The straight
(continous) line is the result obtained using the Ward identities method
\cite{wi,gribov}.}
    \label{fig:zvl}
\end{figure}
\begin{figure}[htbp]
    \centering
    \includegraphics[width=0.6\textwidth]{images/zs.pdf}
    \caption{\itshape $Z_S$ as a function of $\mu^2 a^2$.
We give the renormalization constant obtained by using the
two possible definitions of the quark wave function renormalization,
denoted by $Z_\psi$, eq.\ (\ref{eq:zpsil}), and $Z^\prime_\psi$,
 eq.\ (\ref{eq:zpsi1}).The dashed curve is $Z^{{\rm BSPT}}_S$, from boosted
standard perturbation theory in the infinite volume limit,
 and the full curve is the result obtained in BDPT.}
    \label{fig:zs}
\end{figure}
\begin{figure}[htbp]
    \centering
    \includegraphics[width=0.6\textwidth]{images/za.pdf}
    \caption{\itshape $Z_{A^L}$ as a function of $\mu^2 a^2$.
We give the renormalization constant, obtained by using the
time component $A_0$ and the space components $A_k$, separately.
The dashed line is $Z^{{\rm BSPT}}_A$, from boosted
standard perturbation theory,
 and the full curve is from BDPT. The straight
(continous) line is the result obtained using the Ward identities method
\cite{wi,gribov}.}
    \label{fig:za}
\end{figure}
\begin{figure}[htbp]
    \centering
    \includegraphics[width=0.6\textwidth]{images/zp.pdf}
    \caption{\itshape $Z_P$ as a function of $\mu^2 a^2$. We give the
renormalization constant obtained by using the
two possible definitions of the quark wave function renormalization,
denoted by $Z_\psi$, eq.\ (\ref{eq:zpsil}), and $Z^\prime_\psi$,
 eq.\ (\ref{eq:zpsi1}).The dashed curve is $Z^{{\rm BSPT}}_P$, from boosted
standard perturbation theory in the infinite volume limit,
 and the full curve is the result obtained in BDPT.
}
    \label{fig:zp}
\end{figure}
\begin{figure}[htbp]
    \centering
    \includegraphics[width=0.6\textwidth]{images/zsszp.pdf}
    \caption{\itshape{$Z_S/Z_P$ as a function of $\mu^2 a^2$. The dashed line
corresponds to BSPT, the full line comes from the determination of this
ratio, using the Ward identity method \cite{wi}.
}}
    \label{fig:zsszp}
\end{figure}

In this  section, we give the main results of our numerical study of the
non-perturbative renormalization procedure proposed in this paper and a
comparison
of these results with perturbation theory.

We have performed a
simulation, by generating $36$ independent gluon field configurations,
on a $16^3 \times 32$ lattice, at $\beta=6.0$. The quark propagators have been
computed for a single value of the quark mass ($am_q \sim 0.07$),
 corresponding to a value of the hopping parameter  $K=0.1425$.
All the Green functions have been computed in the lattice Landau
gauge, which  is obtained by
minimizing the functional \beq
{\rm tr}\Bigl[ \sum_{\mu=1}^{4}\Bigl(U_{\mu}(x)+
U^{\dagger}_{\mu}(x)\Bigr)\Bigr].
\eeq
Possible effects from Gribov copies or spurious solutions \cite{giusti} have
not been considered.
As shown below, for those quantities which can be determined also
in a gauge invariant
way, the non-perturbative results obtained
on quark state  in a fixed gauge  are in
good agreement with  those  obtained using the Ward
identity method  \cite{wi}.

\noindent {\bf Vector current:}

In fig.\ \ref{fig:zvl}, the renormalization constant of the local vector current
$Z_{V^L}$, obtained using eqs.\ (\ref{eq:zpsi1}) and (\ref{eq:gvl}), is given.
As expected, $Z_{V^L}$ is independent of the scale,
whithin  statistical errors, up to large values of $\mu^2$, where distortions
due to lattice discretization become important.
This is a consequence of  the
equivalence, which can be established up to terms of $O(\alpha_s a)$, of the
method used here to determine $Z_{V^L}$, and
of the Ward identity for the local vector current, see sec.\ \ref{sec:casivari}.
By using the points at $\mu^2 a^2 \sim 1$, corresponding to $\mu \sim 2$ GeV,
we get $Z^{{\rm NP}}_{V^L}=0.84(1)$, to be compared with
$Z^{{\rm WI}}_{V^L}=0.824(2)$  from the Ward identity method, and
$Z^{{\rm BPT}}_{V^L}=0.83$ ($0.86-0.87$ on a $16^3 \times 32$ lattice,
at $\mu^2 a^2 \sim 1$) from BSPT.  The three
methods are in good agreement.

\noindent {\bf Scalar density.:}

$Z_S$ is an ideal quantity with which to check the validity of
our method for a number of reasons:
it cannot be determined using the Ward identity method, it
is logarithmically divergent and gauge dependent, and
it is not affected by the pole of the Goldstone boson (in contrast to
the pseudoscalar density and the axial current), sec.\ \ref{sec:idea}.

In fig.\ \ref{fig:zs}, $Z_S$ is
shown as a function of the renormalization scale.
We report separately the results, obtained by using
$Z_\psi$, eq.(\ref{eq:zpsil}), and $Z^\prime_\psi$,  eq.(\ref{eq:zpsi1}).
Notice that
the points, corresponding to different definitions of $Z_\psi$, are in very
good agreement where $\mu^2 a^2$ is not too large, and
discretization errors
are small. This can also  be seen  by comparing the results in BSPT
(dashed curve) to those in BDPT (continous curve).
We do not expect agreement
with perturbation theory, either at low $\mu^2$, where higher order effects
become very large, or at high $\mu^2$, where lattice distortions are important,
as shown by the continous curve.
The numerical results follow the theoretical expectations and we find good
agreement between the non-perturbative determination of $Z_S$ and the
predictions
from boosted perturbation theory for $0.3 \le \mu^2 a^2 \le 1$
($1.1$ GeV $\le \mu \le 2$ GeV). Notice the $Z_S$ ($Z_P$) is defined here
from the off-shell Green function and depends on the anomalous dimension
 of the scalar (pseudoscalar) density,
and hence  on log($\mu^2 a^2$). Moreover,
the result is gauge dependent, even in the continuum,
 as can be shown by an explicit
calculation in one-loop perturbation theory.

\noindent {\bf Axial vector current:}

The local axial-vector current shares some of the features of the vector
current.
Its renormalization constant is finite at all orders in perturbation theory
\cite{curci} and can be determined from  Ward identities. However, the
axial-vector current is coupled to the would-be Goldstone boson of $QCD$, which
can give an important contribution at low $\mu^2$.
$Z_{A^L}$, which like $Z_{V^L}$, should be independent of  $\mu$,
 is
shown as a function of $\mu^2$
in fig.\ \ref{fig:za}.

 We interpret the strong $\mu$-dependence of
$Z_{A^L}$, at low $\mu^2$, as the non-\-
per\-tur\-ba\-ti\-ve effect of the pseudoscalar state.
Unfortunately, there is no clear sign of the existence of a plateau between the
non-perturbative regime and the large $\mu$ region, where lattice artefacts
become important\footnote{The plateau could eventually appear more clearly at
larger values of $\beta$.}. Without the information coming
from the Ward identity method \cite{wi}, it would be difficult
to determine $Z_{A^L}$ confidently. Nevertheless,
we do observe that, at  $\mu^2 a^2 \sim 1$, just before lattice artefacts
become large in DPT,  the results are close to
the value determined
with the Ward identity, $Z^{{\rm WI}}_A= 1.06(2)$ \cite{wi,gribov}, see
fig.\ \ref{fig:za}.
As observed in ref.\ \cite{wi},  perturbation theory gives values of
$Z_{A^L}$ smaller than one, while the Ward identity method
gives values larger than one, for $\beta=6.0$--$6.2$. A value
larger than one is also suggested by our non-perturbative
results.

\noindent {\bf Pseudoscalar density:}

The pseudoscalar density
shares the  main features of the scalar density, with two
important differences. Firstly  it is coupled to the would-be Goldstone
boson, and secondly the one-loop
perturbative corrections, with the Clover action, are quite large.
 It is then not
surprising that the non-perturbative value of the corresponding renormalization
constant $Z_P$ lie well below the perturbative result, as shown in fig.
 \ref{fig:zp}.
 Had we used standard perturbation theory in fig.\ \ref{fig:zp},
instead of boosted perturbation theory, the discrepancy would
have been even worse. The discrepancy could have been anticipated from the
results of ref.\ \cite{wi}, where it was shown that the ratio $Z_S/Z_P$,
determined by using the  Ward identity method, was significantly larger than
the result
obtained in boosted perturbation theory. Combining this information
with the agreement between the non-perturbative determination of $Z_S$ and
the result from BPT, one would conclude that the difference is due to
$Z_P$. This may also be due to  the fact that $Z_P$ has
larger one-loop finite corrections, $\sim 35 \%$,
 and that perhaps the  higher order terms, which are not accounted for by
one loop boosted perturbation theory,
are important.

We also present, in fig.\ \ref{fig:zsszp}, $Z_S/Z_P$ as a function of $\mu^2$.
The ratio $Z_S/Z_P$ has a behaviour similar to $Z_{A^L}$. Some discretization
effects appear to be  smaller than those observed in $Z_S$ and $Z_P$
separately.
The numerical results support, qualitatively, the theoretical
interpretation. In the intermediate range of $\mu$, this ratio stays almost
constant.  At small or large values of $\mu$ the results are
unstable, due to non-perturbative or discretization effects respectively.
 At
$\mu^2 a^2 \sim 1 $, the value of this ratio obtained from our
non-perturbative method  is in reasonable agreement with that determined
with the Ward identity method, $Z_S/Z_P=1.64(5)$ \cite{wi}\footnote
{The value published in ref.\ \cite{wi}, using half of the
present stastistics,  was $Z_P/Z_S=0.64(2)$,
corresponding to $Z_S/Z_P=1.56(5)$.}.
